% Figure: the sodium-potassium ATPase as a schematic, with the
% intracellular and extracellular ion concentrations labelled. Shows
% three Na+ out, two K+ in, one ATP hydrolysed.
\begin{figure}[htbp]
\centering
\begin{tikzpicture}[
  membrane/.style={very thick, draw=engred},
  pump/.style={fill=engamber!30, draw=engamber!80!black, very thick,
    rounded corners=2pt},
  na/.style={circle, fill=engblue!70, draw=engblue!90!black,
    inner sep=1.4pt, minimum size=0.36cm, font=\scriptsize,
    text=white},
  k/.style={circle, fill=engred!70, draw=engred!90!black,
    inner sep=1.4pt, minimum size=0.36cm, font=\scriptsize,
    text=white},
  arr/.style={-{Latex[length=2.0mm]}, thick, draw=engblack},
  lbl/.style={font=\scriptsize, align=left},
]

% Plasma membrane (two parallel lines)
\draw[membrane] (-1.0, 1.2) -- (10.0, 1.2);
\draw[membrane] (-1.0, 0.0) -- (10.0, 0.0);

% Pump in the middle
\draw[pump, rounded corners=4pt]
  (4.0, -0.3) -- (5.7, -0.3) -- (5.9, 1.5) -- (3.8, 1.5) -- cycle;
\node[font=\scriptsize\bfseries, align=center] at (4.85, 0.6) {Na$^+$/K$^+$ \\ ATPase};

% Ions on the intracellular side (Na+ inside in)
\node[na] at (3.0, 0.5) {Na};
\node[na] at (3.3, 0.2) {Na};
\node[na] at (2.8, 0.85) {Na};
% ATP
\node[font=\small\bfseries] at (3.1, -0.6) {ATP};

% Ions emerging extracellularly
\node[na] at (6.6, 1.85) {Na};
\node[na] at (7.0, 1.95) {Na};
\node[na] at (7.3, 1.65) {Na};

% K+ on the extracellular side (going in)
\node[k] at (7.0, 1.6) {K};
\node[k] at (6.6, 1.7) {K};

% K+ delivered intracellularly
\node[k] at (4.85, -0.8) {K};
\node[k] at (5.25, -0.65) {K};

% Arrows showing flux direction
\draw[arr] (3.4, 0.5) to[bend left=20]
  node[midway, above, font=\scriptsize] {$3 \times$ out} (5.5, 1.5);
\draw[arr] (4.8, 1.5) to[bend right=20]
  node[midway, below, font=\scriptsize] {$2 \times$ in} (4.4, 0.0);
\draw[arr] (3.4, -0.45) -- (3.95, -0.05)
  node[midway, below right, font=\scriptsize] {hydrolysed};

% Labels for compartments
\node[lbl, anchor=west] at (-0.5, 2.2) {\textbf{extracellular} \\
  $[\text{Na}^+]_o = 145$ mM \\ $[\text{K}^+]_o = 4$ mM};
\node[lbl, anchor=west] at (-0.5, -1.0) {\textbf{cytoplasm} \\
  $[\text{Na}^+]_i = 10$ mM \\ $[\text{K}^+]_i = 140$ mM};

\end{tikzpicture}
\caption[The sodium-potassium ATPase]{The sodium-potassium ATPase
(Na$^+$/K$^+$-ATPase) exchanges three intracellular sodium ions for
two extracellular potassium ions per molecule of ATP hydrolysed
\cite{paper:v6c01-skou-1957}. The net result is a current of one
positive charge out per cycle and the maintenance of the canonical
mammalian ion gradients (Na$^+$: $10\,\text{mM}$ inside,
$145\,\text{mM}$ outside; K$^+$: $140\,\text{mM}$ inside,
$4\,\text{mM}$ outside). The pump is the primary active transporter
that funds the cell's resting potential, the action potential's
restoration phase, and a long list of secondary active transporters
that use the inward sodium gradient as their energy source.}
\label{fig:vol06:ch01:na-k-pump}
\end{figure}
